%%%%%%%%%%%%%%%%%%%%%%%%%%%%%%%%%%%%%%%%%%%%%%%%%
%%%%%%%%%%%%%%%%%%%%%%%%%%%%%%%%%%%%%%%%%%%%%%%%%
%%
%%                       Modèle pour le  Bulletin de l’Association mathématique du Québec
%%
%%                                                              Avril 2013
%%
%%                                       version pour nouveaux utilisateurs de LaTeX
%%
%%%%%%%%%%%%%%%%%%%%%%%%%%%%%%%%%%%%%%%%%%%%%%%%%
%%%%%%%%%%%%%%%%%%%%%%%%%%%%%%%%%%%%%%%%%%%%%%%%%

%%%%%%%%%%%%%%%%%%%%%%%%%%%%%%%%%%%%%%%%%%%%%%%%%
%%%%%%%%%%%%%%%%%%%%%%%%%%%%%%%%%%%%%%%%%%%%%%%%%
%%
%%                     Pour télécharger LaTeX 2e, aller à http://www.ctan.org/tex-archive/CTAN.sites  .
%%                     On arrive à une liste de sites classés géographiquement.
%%		 	Il y a trois sites au Canada: Winnipeg, Waterloo et Halifax.
%%			Par exemple, à Waterloo: http://mirror.csclub.uwaterloo.ca/CTAN/   ,
%%			chercher dans le répertoire "systems"  la version (en jargon: distribution) appropriée
%%			pour votre système d'exploitation...  et suivre les instructions.
%%
%%			Lorsque le logiciel sera téléchargé, ouvrir l'éditeur, par exemple TeXShop,
%%			et depuis le menu déroulant "fichier" choisir  l'option Ouvrir... et choisir le fichier
%%			GabaritAMQ101.
%%			Le fichier GabaritAMQ101.tex qui est un fichier source, s'ouvrira.
%%
%%%%%%%%%%%%%%%%%%%%%%%%%%%%%%%%%%%%%%%%%%%%%%%%%
%%%%%%%%%%%%%%%%%%%%%%%%%%%%%%%%%%%%%%%%%%%%%%%%%

%%%%%%%%%%%%%%%%%%%%%%%%%%%%%%%%%%%%%%%%%%%%%%%%%
%%
%%                                                        AVANT TOUTE CHOSE,
%% 			 enregistrer ce document (GabaritAMQ101.tex) sous un autre nom.
%%                     Par la suite, les modifications s'enregistreront de façon automatique.
%%			 Lire en parallèle le fichier source et la version PDF du document avant d'entrer du texte.
%%			Pour voir le PDF du document final, déclencher la compilation du fichier source
%%			en tapant sur le bouton Typeset (Run typesetting program).
%%
%%============================Votre texte=================================
%% 
%%			Votre texte sera entré ENTRE les commandes \begin{document} et \end{document}.
%%
%%=====================================================================
%%%%%%%%%%%%%%%%%%%%%%%%%%%%%%%%%%%%%%%%%%%%%%%%%

%%%%%%%%%%%%%%%%%%%%%%%%%%%%%%%%%%%%%%%%%%%%%%%%%
%%
%%				Notre Gabarit a été augmenté de plusieurs "packages" au fil de son utilisation 
%%				par nos auteurs potentiels. Aussi, plusieurs commandes que le Gabarit ne
%% 				semble pas documenter sont pourtant disponibles. Vérifier le préambule 
%%				(entre \documentclass[10pt]{article} et \begin{document}) 
%%				avant de faire vos propres ajouts dans la section « Packages propres à l'auteur ».
%%				ou dans la section « Éléments propres à l'auteur ».
%%
%%%%%%%%%%%%%%%%%%%%%%%%%%%%%%%%%%%%%%%%%%%%%%%%%


\documentclass[10pt]{article}

%=====================================Ensemble des packages=====================================									
\usepackage[T1]{fontenc}					% Ce package sert à accueillir les accents et 
									% à faire les césures comme on aime qu'elles soient faites en français.

\usepackage{lmodern}					% Choix de la fonte
\usepackage[utf8]{inputenc}			% Permet de traiter les textes accentués sur un Mac.
									% Les utilisateurs de Unix doivent remplacer [applemac] par [latin1],
									% et ceux de Windows doivent remplacer [applemac] par [ansinew].
									% Remplacer [applemac] par [utf8] permet d'utiliser l'encodage utf-8 (Unicode) 																					% censé être universel, mais qui est parfois mal supporté par les systèmes
									% plus anciens.
									% Ensuite, entrer le texte accentué comme à l'ordinaire.
\usepackage[english,french]{babel} 			% Permet de rendre le fichier maître plus compatible,
									% en particulier avec l'insertion de figures,
									% (sauf celle de figures flottantes (package floatingfigure))
									% et avec la définition des caractères à doubles jambes.
\usepackage{graphicx}
									% Le fichier maître est aussi compatible avec l'utilisation de listes avec 																						% \itemize et avec l'utilisation de notes de bas de page.
									
									% Le fichier maître dans sa version du 21 novembre 2012 est 																								% compatible avec les packages hyperref et url.
									% Il faudrait vérifier s'il l'est avec  \ref{}, \eqref{}.
									 
\usepackage{amssymb,latexsym,amsmath}	% Packages qui permettent l'utilisation des symboles et des commandes
									% mathématiques de AMS-LaTeX,
									% et en particulier \align qui pour l'instant ne fonctionne pas.
\usepackage{mathrsfs}					% Package qui donne accès à une police calligraphique, comme eucal.
\usepackage{multirow}					% Package qui permet de fusionner des lignes ou des colonnes de tableaux.
\usepackage{euscript}					% Package qui donne accès à la fonte Euler.
\usepackage{fancyhdr}					% Package qui permet de définir ses propres en-têtes et pieds de page.
\usepackage{pdfsync}					%  Package qui permet de passer du texte sur le fichier
									% pdf à la ligne du fichier source correspondant.
\usepackage{theorem}					% Package permettant de définir son propre style
									% de présentation des théorèmes.
\usepackage{wrapfig}					% Package qui permet au texte de s'enrouler autour de flotants. 
\usepackage{floatflt}						% Package qui permet au texte de s'enrouler autour de flotants. 
\usepackage[protrusion=true]{microtype}          %Package qui "rectifie" les caractères qui dépassent dans la marge.

\usepackage[all]{xy}						%Permet de tracer des graphes fléchés.
									%Voir http://www.tug.org/TUGboat/tb22-4/tb72perlax.pdf

\usepackage{geometry}					%Permet de définir les dimensions de la partie imprimée de la page.
									%Voir http://ctan.cms.math.ca/tex-archive/macros/latex/contrib/geometry/geometry.pdf
\usepackage{tikz}						%Permet, entre autres, de mettre du texte en filigrane.
									%Voir http://math.et.info.free.fr/TikZ/bdd/TikZ-Impatient.pdf 
									%p. 165
												
\usepackage{booktabs}					% Permet d'utiliser le gabarit de tableaux utilisant les commandes
									% of \toprule, \midrule and \bottomrule pour les lignes horizontales dans les tableaux. 
\usepackage{color}
\usepackage{url}						%Permet la présentation d'adresses url sous leur forme habituelle.
									%Voir http://www.tex.ac.uk/cgi-bin/texfaq2html?label=setURL
\usepackage[colorlinks=true,urlcolor=blue, citecolor=cyan, linkcolor=magenta]{hyperref}	
									%Permet de rendre actifs les liens créés avec le package précédent.
									%Voir /usr/local/share/texmf-dist/doc/latex/hyperref/manual.pdf??????????????????????
									%Pour une adresse de courriel: \href{mailto:foo@bar.abc}{Good Link!}
\usepackage{esvect}						%Pour de belles flèches sur les vecteurs, cadeau de Jean-Philippe Morin
%==================================Fin de l'ensemble des packages==================================


%==================================Pour de belles flèches sur les vecteurs, cadeau de Jean-Philippe Morin=====
\renewcommand{\vec}[1]   {\protect\vv{#1}}     
                          % notation de vecteur (\vv définie par package esvect)
                          % Exemple: \vv{v} ou \vv{AB} ou \vv{A_1A_2}
%%============================================================================================ 


%%%====================== Début macro  pour numérotation des équations selon les sections du texte==========
%\makeatletter
%\renewcommand\theequation{\thesection.\arabic{equation}}
%\@addtoreset{equation}{section}
%\makeatother
%%%======================  macro: mode d'emploi ==================================================
%%%    			Pour activer cette macro commande, enlever le symbole % au début de la ligne suivante
%\numberwithin{equation}{section} 

%%============================================================================================ 


%%=============================Les accents===============================
%%   La commande  usepackage[applemac]{inputenc} sert à traiter les textes accentués.
%%  Les utilisateurs de unix ou d'un PC sous Windows doivent remplacer [appelmac] par [latin1].
%%  Il existe d'autres options de codage pour d'autres systèmes d'exploitation.
%%  Voir à ce sujet http://www.grappa.univ-lille3.fr/FAQ-LaTeX/11.3.html.

%%=====================================================================


%%==========================Résumé: mode d'emploi==========================
%%
%%    Entrer le résumé du texte entre les commandes \begin{abstract} et \end{abstract} plus bas.
%%
%%=====================================================================



%%===========Les énoncés ayant la même présentation que les théorèmes==================================
\newtheorem{theorem}{Théorème}
%[section]
\newtheorem{lemma}{Lemme}
%[section]
\newtheorem{proposition}{Proposition}
%[section]
\newtheorem{corollary}{Corollaire}
%[section]
\newtheorem{definition}{Définition}
%[section]



\newtheorem{nota}{Notation}
%[section]
\newtheorem{ex}{Exemple}
%[section]
\newtheorem{nb}{N.B.}
%[section]
\newtheorem{remark}{Remarque}
%%============================================================================================


%%===================Théorèmes: mode d'emploi=============================
%% Commencer l'énoncé d'un théorème par \begin{theorem} et la terminer par \end{theorem}.
%% Commencer l'énoncé d'un lemme par \begin{lemma} et la terminer par \end{lemma},etc.
%%=====================================================================


%%==========================Les démonstrations============================
\newenvironment{proof}
{\par\noindent
\textit{Démonstration}\ }
{
\hfill{$\Box$}
}
%%====================== Démonstrations: mode d'emploi======================
% %       Commencer une démonstration par  \begin{proof} et la terminer par \end{proof}.
%%=====================================================================


%%============================Guillemets français===============================================
%%        Pour obtenir les guillemets français:
%%        \og pour ouvrir les guillemets
% %       \fg pour fermer les guillemets
%%==========================================================================================

%%==============================Les siècles====================================================
%%        Pour la notation des siècles en chiffres romains
%%        La commande est  \siecle{#}
%%        Exemple  \siecle{21}.
%%==========================================================================================

\newcommand{\siecle}[1]{%
\ifnum #1=1%
\uppercase\expandafter{\romannumeral #1}\textsuperscript{er}%
\else%
\uppercase\expandafter{\romannumeral #1}\up{e}%
\fi}

%%==============================Les angles mesurés en degrés====================================
%%        Pour la notation des mesures d'angles en degrés, il y a deux possibilités:
%%        En mode math, la commande est  #^{\circ}
%%        Exemple  45^{\circ} donne 45 degrés.
%%
%%        En mode texte, la commande est  \textdegree
%%        Exemple  45\textdegree  donne 45 degrés.%===========================================================================================

%%==============================Les adresses url===============================================
%%        Pour écrire une adresse url de façon claire et la rendre active, on utilise le package
%%        hyperref (qui doit être le dernier invoqué) et la commande \href:
%%	      La commande est  \href{une adresse url}{un texte d'appel qui s'imprimera}
%%	      Si vous souhaitez que l'adresse url apparaisse dans le PDF, 
%%	      alors utiliser  \href{une adresse url}{la même adresse}
%%	      Si jamais il y a un caractère spécial (et donc interdit dans du texte) dans l'adresse, 
%%	      alors placer un « \ » devant ce caractère dans la seconde paire d'accolades , mais ne changez rien dans la première 
%%        Exemple1  \href{http://newton.mat.ulaval.ca/amq/}{http://newton.mat.ulaval.ca/amq/}
%%        Exemple2  \href{http://mathdoc.emath.fr/archives-bourbaki/}{Archives de Bourbaki}
%%        Vous trouverez des exemples d'utilisation dans l'identification de l'auteur du texte
%%        et dans la bibliographie.
%% 
%%==========================================================================================

%%=====================================================================
%                                   D\'{e}finition des marges pour le Bulletin (NE PAS MODIFIER)
\geometry{height=18.6cm,width=14.4cm}
\geometry{hmargin={3.6cm,3.6cm}}
\geometry{vmargin={4.3cm,5cm}}
%%Ces trois lignes doivent rester avant le \begin{document}.
\parindent=0.0in
\parskip=0.1in

\setcounter{page}{1}
%%=====================================================================

%%========================Caractères particuliers usuels======================
%%
%%==============================Double jambe.============================
\DeclareSymbolFont{AMSb}{U}{msb}{m}{n}
\DeclareSymbolFontAlphabet{\Bbb}{AMSb}
\newcommand{\ds}{\displaystyle}

%%=========================Double jambe: mode d'emploi=====================
%%        La commande est : {\Bbb}.
%%        Exemple $\Bbb{N}$ désigne l'ensemble des entiers naturels.
%%=====================================================================

%%==================numérotation des lignes du PDF==========================
%%        Pour que les lignes du document PDF généré par LaTeX soient numérotées,
%%        on utilise le package lineno disponible sur
%%        ftp://ftp.inria.fr/pub/TeX/CTAN/macros/latex/contrib/lineno/
%%        ainsi que les 17 lignes de code situées entre \usepackage[]{lineno}
%%        et \usepackage{url}.
%%        Toutefois, dans l'état actuel du gabarit la numérotation n'est pas activée.
%%         Pour activer cette fonction, enlever les symboles % situés devant 
%%        \begin{linenumbers} et \end{linenumbers}.
%%=====================================================================       

\usepackage[]{lineno}

\newcommand*\patchAmsMathEnvironmentForLineno[1]{%
  \expandafter\let\csname old#1\expandafter\endcsname\csname #1\endcsname
  \expandafter\let\csname oldend#1\expandafter\endcsname\csname end#1\endcsname
  \renewenvironment{#1}%
     {\linenomath\csname old#1\endcsname}%
     {\csname oldend#1\endcsname\endlinenomath}}% 
\newcommand*\patchBothAmsMathEnvironmentsForLineno[1]{%
  \patchAmsMathEnvironmentForLineno{#1}%
  \patchAmsMathEnvironmentForLineno{#1*}}%
\AtBeginDocument{%
\patchBothAmsMathEnvironmentsForLineno{equation}%
\patchBothAmsMathEnvironmentsForLineno{align}%
\patchBothAmsMathEnvironmentsForLineno{flalign}%
\patchBothAmsMathEnvironmentsForLineno{alignat}%
\patchBothAmsMathEnvironmentsForLineno{gather}%
\patchBothAmsMathEnvironmentsForLineno{multline}%
}


%%==========================Packages propres à l'auteur==============================
%%
%%
%%==========================Fin des packages propres à l'auteur========================

\usepackage{url}

  \usepackage{hyperref}
 
 %%============================pour des barres de fraction obliques=======================
 
%running fraction with slash - requires math mode.
\newcommand*\rfrac[2]{{}^{#1}\!/_{#2}}

%		\rfrac{3}{7}
%%===============================================================================
  
  
  

%%================pour des barres de fraction latérales, une petite , une grande et une très grande===============
\newcommand{\fracinline}[2]{\raisebox{0.4ex}{$#1$} / \raisebox{-0.7ex}{$#2$}}
\newcommand{\bigfracinline}[2]{\raisebox{0.8ex}{$#1$}  \Big/ \raisebox{-1.4ex}{$#2$}}
\newcommand{\biggfracinline}[2]{\raisebox{0.8ex}{$#1$}  \Bigg/ \raisebox{-1.4ex}{$#2$}}

%%		\fracinline{1 + \cos x}{\sin x}
%%		\bigfracinline{1 + \cos x}{\sin x}
%%		\biggfracinline{1 + \cos x}{\sin x}
%%=======================================================================================



%%=====================================================================
%%                                                              Fin du préambule
%%=====================================================================
\everymath{\displaystyle}
\begin{document}


%%==========================Éléments propres à l'auteur==============================
%%
%%
%%==========================Fin des éléments propres à l'auteur========================


%%===========================Pour mettre le titre des tableaux en français===============================
%%
 %\def\tablename{{\scshape Tableau}}%
 \renewcommand{\tablename}{\textsc{Tableau}}
%%
%%=========================Fin de la nouvelle définition=============================================

%\begin{linenumbers}

%========================Pour mettre "du texte " en filigrane sur la page===========
%
%	 \begin{tikzpicture}[remember picture,overlay]
%	\node[rotate=60,scale=10,text opacity=0.1] at (current page.center) {du texte};
%	\end{tikzpicture}
%======================================================================

\begin{center}
\textsf{\LARGE \textbf{Titre du texte, si possible moins de 50 caractères}}
\end{center}

\begin{flushright}
\textsc{Joseph Bleau, département de mathématiques,\\
C\'{e}gep de Fleurimont}\\
\href{mailto:jo-bleau@cegepfleurimont.qc.ca}{contact par courriel avec l'auteur}\\
  \href{http://archimede.mat.ulaval.ca/amq/}{le site de l'AMQ}\\
\end{flushright}

\baselineskip=1.2\baselineskip

%%=====================================================================
%%          Résumé
%%=====================================================================

\begin{abstract}

Dans ce Gabarit, on devra entrer tous les textes en caractères accentués, comme à l'habitude\footnote{Si vous utilisez les accents habituels, vous nous rendez le grand service de nous éviter d’avoir à modifier un texte qui ne les utiliserait pas.}.\\
 Le résumé apparaît ici. Long d'au maximum 100 mots, il peut contenir des expressions en  {\em italique}, en {\bf caractères gras}, entre \og guillemets\fg~ou en  \textsc {petites capitales}. On peut aussi utiliser des symboles particuliers comme $\Bbb{R}$ ,  $\mathscr{F}$ ou encore $\mathscr{P}(E)$ et même des expressions mathématiques comme $f_k(x)=a x^2+bx+c$, ou encore  «  l’ensemble $\mathfrak{F}$ des formules du système formel $\mathfrak{L}$ est bien défini ».

Il est possible de reculer dans le temps. Ainsi, on peut parler du  \siecle{1} siècle de notre ère ou des  \siecle{20} et  \siecle{21} siècles.


\end{abstract}
%%=====================================================================
%%          Fin du résumé
%%=====================================================================


\bigskip



%%===============================Section 1===============================
\section{Structuration du texte}\label{sec1}
\bigskip
%%==============================Définition 1===============================
\bigskip
\begin{definition}
On appelle section toute partie du texte incluse entre une annonce de section et la suivante. La même convention s'applique aux sous-sections.
\end{definition}

 
Les sections et sous-sections sont numérotées de façon automatique. La même règle s'applique, catégorie par catégorie, aux théorèmes, aux lemmes, aux corollaires et aux définitions.

Toutefois, on trouvera un moyen de remplacer la numérotation catégorie par catégorie par une numérotation en fonction de la section ou de la sous-section dans le préambule et c'est là aussi que vous devrez l'activer. Vous trouverez plus d'information sur ce sujet et sur bien d'autres, dans M.-P. Kluth et B. Bayart, {\em FAQ  \LaTeX~de l'équipe Grappa} \cite{faq}.

On trouvera une liste de listes de symboles disponibles pour écrire des expressions mathématiques pages 50 à 56, dans T. Oetiker, H. Partl, I. Hyna et E. Schlegl (1999). {\em Une courte (?) introduction à \LaTeXe, ou  \LaTeXe~en 88 minutes}  \cite{oeti}.



%%===============================Section 2===============================
\section{Les expressions mathématiques et les équations}\label{sec2}

%%============================Sous-section 2.1============================
\subsection{Expressions mathématiques au fil du texte}\label{fil} 
Le signe de dollar permet d'insérer au fil du texte des expressions mathématiques qui apparaîtront en italique. Ainsi, lorsqu'on utilise une fonction f, pour qu'elle apparaisse comme un objet mathématique, il suffit de l'écrire $f$, entre deux signes de dollar. Les expressions peuvent être plus complexes:\\ Soit, $f: \Bbb{N} ^{2*} \rightarrow \Bbb{R}$ telle que pour tout $(m, n) {\in} \Bbb{N} ^{2*} $ on a $f(m+1,n+1) = f(m+1,n) + f(m,n+1)$. 


%%============================Sous-section 2.2============================
\subsection{Expressions mathématiques sur une ligne dédiée}\label{dediee} 
On peut souhaiter que l'équation soit mise en évidence sur une ligne qui lui soit dédiée:\\

Soit, $f: \Bbb{N} ^{2*} \rightarrow \Bbb{R}$ telle que pour tout $(m, n) {\in} \Bbb{N} ^{2*} $ on a
\begin{equation*}
f(m+1,n+1) = f(m+1,n) \qquad+ f(m,n+1).
\end{equation*}

On peut souhaiter qu'en plus cette équation soit numérotée et qu'on puisse faire référence plus loin dans le texte à ce numéro:\\
Soit, $f: \Bbb{N} ^{2*} \rightarrow \Bbb{R}$ telle que pour tout $(m, n) {\in} \Bbb{N} ^{2*} $ on a
\begin{equation}
f(m+1,n+1) = f(m+1,n) + f(m,n+1). \label{recur}
\end{equation}

Pour faire référence à une équation numérotée, on utilise l'étiquette qui lui a été attribuée. Ainsi on peut invoquer l'équation (\ref{recur}) de la section (\ref{dediee}). Ce système est résistant aux changements de place dans le texte. Il est possible qu'il soit nécessaire de compiler plusieurs fois après une modification du texte pour assurer le suivi des références.

%%============================Sous-section 2.3============================
\subsection{Expressions mathématiques sur plusieurs lignes, avec numérotation }\label{numero} 

$AC=CB$ et $AB=AC$
Il arrive qu'on ait une suite d'égalités à présenter et qu'on souhaite aligner les signes d'égalité.
On traite alors la suite d'équations comme un tableau à 3 colonnes délimitées  par le symbole de l'esperluette (et commercial).
\newpage
Exemple:

\begin{theorem} Si $\vv{v} \neq \vv{0}$, alors le vecteur $\frac{1}{\| \vv{v}\|}\vv{v}$, que l'on écrit parfois $\frac{\vv{v}}{\| \vv{v}\|}$, est un vecteur unitaire de même direction et de même sens que $\vv{v}$.
\end{theorem}

La partie du théorème précédent portant sur la norme se démontre ainsi:

\begin{proof}
\begin{eqnarray}
\Big\|{\frac{1}{\| \vv{v}\|}\vv{v}} \Big\|&=& \Big |\frac{1}{\| \vv{v}\|} \Big |\| \vv{v}\|\label{norme}\\
&=&\frac{1}{\| \vv{v}\|}\| \vv{v}\|  \nonumber\\
&=&\frac{\| \vv{v}\|}{\| \vv{v}\|} \nonumber\\
&=&1 \label{nombre}
\end{eqnarray}
La comparaison des équations (\ref{norme}) et (\ref{nombre}) démontre le théorème.\end{proof}


Pour obtenir une suite d'équations non numérotées, 

\begin{eqnarray*}
\Big\|{\frac{1}{\| \vv{v}\|}\vv{v}} \Big\|&=& \Big |\frac{1}{\| \vv{v}\|} \Big |\| \vv{v}\| \\
&=&\frac{1}{\| \vv{v}\|}\| \vv{v}\| \\
&=&\frac{\| \vv{v}\|}{\| \vv{v}\|} \\
&=&1
\end{eqnarray*}

%%============================Sous-section 2.4============================
\subsection{Expressions mathématiques d'une fonction définie par morceaux }\label{morceaux} 

Un classique:
%%==============================Définition 2==============================
\begin{definition}
On appelle valeur absolue la fonction
{f}: $\Bbb{R} \rightarrow \Bbb{R}_+$ définie par
\[f(x)=
\begin{cases}
\;\;x & \text{, si } x \geq 0 \\
-x & \text{, sinon}.
\end{cases}
\]
\end{definition}

ou encore 
\[f(x)=
\begin{cases}
\; \;x & \text{, si } x > 0, \\
\; \;0& \text{, si } x = 0, \\
-x & \text{, \emph{sinon}}.
\end{cases}
\]

%%===============================Section 3===============================
\section{Les tableaux}\label{tab}

\bigskip

Ce tableau (\ref{Tableau1}) comporte 2 colonnes, dans lesquelles le texte sera centré, et 3 lignes.
Sur la première ligne, les deux colonnes sont fusionnées.

 \begin{table}[ht]
\begin{center}
\begin{tabular}{|c|c|}
\hline
\multicolumn{2}{|c|}{Tableau}\\
\hline
Cellule 1&Cellule 2\\
\hline
Cellule 3&Cellule 4\\
\hline
\end{tabular}
\end{center}
\caption{ \label{Tableau1} {Mod\`{e}le de tableau dans l'environnement table} }
\end{table}

On peut l'invoquer ailleurs dans le texte car il porte une étiquette. Ainsi, le tableau (\ref{Tableau1}) ne comporte que des lignes simples, alors que dans le tableau (\ref{Tableau2}) on a inséré des doubles lignes horizontales et des doubles lignes verticales.

 \begin{table}[ht]
\begin{center}
\begin{tabular}{|l||r| c| c| c| c|| c c|}
 \hline
 Nom &  Pr\'{e}nom & Institution & Article &  & &  &  \\ \hline \hline
   Fayard & Clotilde& École maternelle de la Métare& La maquette & 3 & 4& 5 & 6 \\ \hline
   1 &1& 2 & 3 & 4& 5& 6 &  0 \\ \hline
   1 &1& 2 & 3 & 4& 5& 6 &  0 \\ \hline
   1 &1& 2 & 3 & 4& 5& 6 &  0 \\ \hline
   1 &1& 2 & 3 & 4& 5& 6 &  0 \\ \hline

\end{tabular}
\end{center}
\caption{ \label{Tableau2} {Autre mod\`{e}le de tableau dans l'environnement table} }
\end{table}


Les tableaux, comme les figures, sont des objets flottants. Cela signifie que le logiciel les insère immédiatement s'il y a suffisamment de place pour qu'il tienne dans la page. Sinon, le logiciel remplit la page avec du texte et insère le tableau en haut de la page suivante.




 \begin{table}[ht]
 \begin{flushright}
\begin{tabular}{ll}\ 
Estragon.   & -  […] Allons-nous-en. \\
~Vladimir.   & -  On ne peut pas. \\
~Estragon.   & -  Pourquoi ? \\
~Vladimir. & -  On attend Godot. \\ 
 \multicolumn{2}{r}{~~~~~~Samuel Beckett, \begin{itshape} En attendant Godot \end{itshape}} \\
 \end{tabular}
 \end{flushright}
\caption{ \label{Tableau3} {Tableau dramatique} }
\end{table}

\textsc{Attention}, cet autre tableau\ref{tab:template} est un peu plus sophistiqué. Lire les instructions dans le fichier source.
\bigskip

\begin{table} 									% Pour préciser la position souhaitée pour le tableau, 
											% ajouter immédiatement après la parenthèse fermante
											%  soit  [h] pour  ici, [t] pour en haut, [b] pour en bas, ou [p] pour sur une page séparée 
\centering 										% Le tableau sera centré dans la page, comment out to left-justify
\begin{tabular}{l c c c c c} 							% La dernière paire de parenthèses définit d'une part le nombre de colonnes,
											%  d'autre part la position des séparateurs de colonnes (barres verticales (|)),
											% et enfin la position du texte à l'intérieur de la colonne 
											% (c pour au centre, l pour à gauche et r pour à droite).
											% Pour spécifier la largeur, utiliser p{largeur}, par exemple p{5cm}.
\toprule 										% Trace une ligne horizontale au dessus du tableau.
& \multicolumn{5}{c}{Milieu de croissance} \\ 			% On peut fusionner plusieurs colonnes en une seule cellule en utilisant la commande
											%  \multicolumn comme sur cette ligne
\cmidrule(l){2-6} 								% Trace une ligne horizontale sur une partie de largeur du tableau,
											% En ajoutant  (r) ou (l) juste avant la paire de parenthèses, 
											% on indique s'il faut raccourcir la ligne du côté gauche (l) ou du côté droit (r).
Variété & 1 & 2 & 3 & 4 & 5\\ 						% Ceci est la ligne des titres de colonnes. 
\midrule 										% Ligne horizontale à l'intérieur du tableau.
GDS1002 & 0.962 & 0.821 & 0.356 & 0.682 & 0.801\\ 	% Contenu de la ligne no 1
NWN652 & 0.981 & 0.891 & 0.527 & 0.574 & 0.984\\ 	% Contenu de la ligne no 2
PPD234 & 0.915 & 0.936 & 0.491 & 0.276 & 0.965\\ 		% Contenu de la ligne no 3
JSB126 & 0.828 & 0.827 & 0.528 & 0.518 & 0.926\\		% Contenu de la ligne no 4
JSB724 & 0.916 & 0.933 & 0.482 & 0.644 & 0.937\\ 		% Contenu de la ligne no 5
\midrule										% Ligne horizontale à l'intérieur du tableau.
\midrule 										% Ligne horizontale à l'intérieur du tableau.
Taux moyen & 0.920 & 0.882 & 0.477 & 0.539 & 0.923\\	 % Summary/total row
\bottomrule 									% Trace une ligne horizontale en dessous du tableau.
\end{tabular}
\caption{Texte de la légende}						% Texte de la légende. Si la légende n'est pas nécessaire, on peut commenter le tableau ailleurs 																	% dans le texte.
\label{tab:template} 								% Étiquette qui permet d'invoquer le tableau ailleurs dans le texte,
											% Les références dans le texte seront de la forme \ref{étiquette}, par exemple \ref{tab:template}.
\end{table}


%%===============================Section 4===============================
\section{Les figures}\label{fig}

On insère les figures en formats JPEG ou PDF.

%%============================Sous-section 4.1============================
\subsection{Insertion d'une figure avec légende étiquetée}\label{fig+eti} 

\begin{figure}[ht]
\centerline{\includegraphics[scale=.8]{pics/Figure1}} \caption{ \label{Modele figure1} Mod\`{e}le d'une figure enregistrée en format JPEG avec  légende étiquetée }
\end{figure}

\begin{figure}[ht]
\centerline{\includegraphics[scale=.99]{pics/Figure2}} \caption{ \label{Modele figure2} Mod\`{e}le d'une figure enregistrée en format PDF avec  légende étiquetée }
\end{figure}

Comme pour les équations numérotées ou les sections, on peut invoquer une figure dont la légende a été étiquetée sans craindre les modifications au texte en cours de composition. 

Ainsi, on peut invoquer les figures (\ref{Modele figure1}) et (\ref{Modele figure2}) pour attirer l'attention du lecteur sur la manière de contrôler la taille d'une figure.

%%============================Sous-section 4.2============================
\subsection{Insertion d'une figure avec légende manuelle}\label{leg+manu} 

\begin{center}
\includegraphics[scale=.99]{pics/Figure1}\\
{\small \textsc{Figure 3} -- Mod\`{e}le d'une figure avec l\'{e}gende manuelle}
\end{center}

On ne peut pas invoquer cette figure autrement que manuellement. Nous n'encourageons pas l'insertion de figures de cette façon car dans ce cas il arrive que la figure elle-même et son titre se retrouvent sur deux pages différentes.

Tout ce qui a été étiqueté, sections, sous-sections, figures, tableaux et entrées de la bibliographie, peut être invoqué ailleurs dans le texte.


%%===============================Section 5===============================
\section{Les citations et les notes de bas de page}\label{sec5}

On peut faire des citations. Les citation courtes (5 lignes ou moins) apparaîtront au fil du texte, entre guillemets français \og Diophante vivait à une époque où les mathématiques alexandrines perdaient leur puissance créatrice.\fg~disent Dahan-Dalmedico et Peifer\cite{daha}, p. 65.

Une citation longue (plus de 5 lignes) apparaîtra dans un paragraphe (ou plus selon le texte cité) qui lui sera dédié.\\ 
 \begin{quote}  Avec Diophante, un nouveau chapitre des mathématiques s'ouvre et il est impossible de mettre en lumière le courant dont il est l'aboutissement. La vie de Diophante est très peu connue, et la période précise pendant laquelle il a vécu reste contestée (\siecle{3} siècle après J.-C.). Sa grande oeuvre, {\ les Arithmétiques}, devait comprendre, d'après ce qu'il écrit lui-même dans l'introduction, treize livres.
 
 Depuis le \siecle{16} siècle, seuls six livres étaient connus. Ils provenaient d'un manuscrit grec découvert en 1464 par Regiomontanus à Venise, qui était la copie d'un manuscrit plus ancien.(Dahan-Dalmedico et Peifer\cite{daha}, p. 72)\end{quote} 
 
Il est aussi possible d'insérer des notes de bas de page. Toutefois, celles-ci sont à utiliser avec la plus grande parcimonie, et seulement quand l'auteur pense vraiment qu'il lui est impossible d'inclure une incise dans son texte \footnote {Un texte de 10 pages avec 12 notes de bas de page n'est pas acceptable.}. Les références bibliographiques, notamment, ne doivent jamais figurer en notes, mais doivent être incluses dans le texte sous les normes données ci-après. Toutes les références correspondant à celles données dans le texte doivent figurer dans la bibliographie. Nous attirons l'attention des auteurs sur le fait que l'exactitude en matière de références bibliographiques requiert que soit entreprise une vérification, voire une recherche documentaire qui est de la responsabilité de l'auteur. 



%%===============================Section 6===============================
\section{Conclusion}\label{sec6}

C'est à la fin de cette section qu'on retrouvera le cas échéant de brefs remerciements.

{\em Remerciements} 
à Leslie Lamport\cite{lamp}, pour la qualité impressionnante de son travail et ses précieux conseils pour la rédaction de ce document, et aux Annales des sciences mathématiques du Québec dont le gabarit nous a inspirés.

%%===============================Section 7===============================
%%===================R\'{e}f\'{e}rences bibliographiques MANUELLES=============
%%
%%
%%                                                           NE PAS UTILISER BIBTEX
%%
%%       La liste des références apparaîtra entre les commandes
%%       \begin{thebibliography}{10} et \end{thebibliography}{10},
%%       par ordre alphabétique du premier auteur.
%%
%%       Toutes les références bibliographiques apparaissant dans la liste des références
%%        devraient être citées dans le corps du texte. Réciproquement, toutes les références
%%        citées dans le corps du texte devraient apparaître en fin de texte dans la liste des 
%%        références.
%%
%%       Chaque item de cette liste sera entré grâce à la commande \bibitem.
%%
%%       Pour un livre:
%%        \bibitem{étiquette} Nom de l'auteur, Initiale. (année de publication). {\em Titre du livre} (Numéro d'édition éd.) (vol. numéro de volume) (Nom du traducteur trad.). Ville, Pays (sauf pour les USA où on mettra Ville, abr. postale de l'état): Éditeur.
%%
%%       Pour un article publié dans une revue:
%%       \bibitem{étiquette} Nom de l'auteur, Initiale. (année de publication). Titre de l'article, {\em Nom de la revue, Numéro du volume} (no. numéro de fascicule), numéro de la première page--numéro de la dernière page.
%%
%%       Pour un chapitre dans un ouvrage collectif (plusieurs auteurs):
%%       \bibitem{étiquette} Nom de l'auteur, Initiale. (année de publication). Titre du chapitre, Dans Initiale. Nom du directeur (dir.) {\em Titre de l'ouvrage}, (numéro d'édition ou de chapitre, p. numéro de la première page--numéro de la dernière page), Ville, Pays: Éditeur.
%%
%%       Pour un document en lige
 %%       \bibitem{étiquette} Nom de l'auteur, Initiale. (date de création, ou année de création mise à jour date de mise à jour). Titre du document (pour une section d'un document long, traiter comme pour un chapitre de livre en remplaçant les numéros de pages par le numéro de chapitre ou de section), **.
%%
%%       Pour invoquer une référence dans le texte, utiliser la commande
%%        \cite{étiquette}
%%       
%%        *Pour plus de renseignements, consulter 
%%        http://benhur.teluq.uqam.ca/~mcouture/apa/normes_apa_francais.pdf
%%        
%%        ** Pour les informations de localisation, consulter le même document
%%        
%%====================================================================


\bigskip

\bibliographystyle{alpha}
\begin{thebibliography}{10}

\bibitem{cout} Couture, M. (2010, mise à jour 18 avril). {\em Revue internationale des technologies en pédagogie universitaire. Normes bibliographiques - Adaptation française des normes de l’APA}. Récupéré le 3 janvier 2011 du site de l’auteur :  \href {http://www.teluq.uqam.ca/~mcouture/apa}{http://www.teluq.uqam.ca/~mcouture/apa}

\bibitem{daha} Dahan-Dalmedico, A.  et Peiffer, J. (1982). {\em Histoire des mathématiques-Routes et dédales}. Paris, France: Études vivantes.

\bibitem{desg} Desgraupes, B. (2003). {\em \LaTeX, Apprentissage, guide et référence}. France: Vuibert.

\bibitem{ctan} Fairbairns, R. (2006, mise à jour le 21 janvier 2006). {\em Search CTAN For a File}, Récupéré le 3 janvier 2011 du site de CTAN :\href{http://www.tex.ac.uk/CTANfind.html}{http://www.tex.ac.uk/CTANfind.html}.\\
Demander texlive pour arriver à \\tex-archive/systems/texlive/Images/Images/texlive-inst.iso.bz2

Cliquer pour récupérer l'image compressée du CD au format bzip2 (550 Mo!)

\bibitem{flipo} Flipo, D. (2004).  \LaTeX, logiciel libre et gratuit, comme alternative à MS-Word, {\em Quadrature}, no. 54, 12--15. 

\bibitem{faq}  Kluth, M.-P. et Bayart, B. (1997, mise à jour 16 octobre 2001). {\em FAQ  \LaTeX~de l'équipe Grappa}. Récupéré le 3 janvier 2011 du site du groupe Grappa de l'Université de Lille 3 : \href{http://www.grappa.univ-lille3.fr/FAQ-LaTeX/}{http://www.grappa.univ-lille3.fr/FAQ-LaTeX/}.

\bibitem{lamp} Lamport, L. (1986). {\em \LaTeX: A Document Preparation System}. Reading, MA: Addison-Wesley Publishing Company.

\bibitem{mass} Masson, B. (2009). \begin{itshape}{\LaTeX, créer ses commandes}\end{itshape}. Récupéré le 3 janvier 2011 du site de l'auteur :\href {http://bertrandmasson.free.fr/wp-content/uploads/commande.pdf}{http://bertrandmasson.free.fr/wp-content/uploads/commande.pdf}.

\bibitem{oeti} Oetiker, T., Partl, H., Hyna, I. et Schlegl, E. (1999). {\em Une courte (?) introduction à \LaTeXe, ou  \LaTeXe~en 88 minutes}, Version 3.3. Récupéré le 3 janvier 2011 du site du Loria : \href{http://tex.loria.fr/general/flshort-3.3.pdf}{http://tex.loria.fr/general/flshort-3.3.pdf}. 

\bibitem{roll}Rolland, C. (2000). {\em \LaTeX~par la pratique} ($2^e$ éd.). Paris, France: Eyrolles.

\end{thebibliography}

%\end{linenumbers}

\end{document}

%%%%%%%%%%%%%%%%%%%%%%%%%%%%%%%%%%%%%%%%%%%%%%%%%
%%%%%%%%%%%%%%%%%%%%%%%%%%%%%%%%%%%%%%%%%%%%%%%%%
%%
%%                       Les auteurs:
%%
%%                                    Marie-Claude Côté                 
%%                                    Fernand Beaudet
%%                                    Marie-Jane Haguel
%%
%%
%%                                                                                                                          version novembre 2010
%%%%%%%%%%%%%%%%%%%%%%%%%%%%%%%%%%%%%%%%%%%%%%%%%
%%%%%%%%%%%%%%%%%%%%%%%%%%%%%%%%%%%%%%%%%%%%%%%%%

